\setcounter{section}{1}

\Needspace{5\baselineskip}
\hypertarget{vapo}{%
\section{VAPO}\label{vapo}}

\begin{quote}
This article is a paper-reading note. Its content represents the corresponding paper's method or the author's understanding and should not be treated directly as field consensus or engineering best practice.
\end{quote}

Value-model-based Augmented Proximal Policy Optimization (VAPO) returns to actor--critic. It is an improved PPO algorithm that also incorporates design ideas from PPO and GRPO/DAPO.

\Needspace{5\baselineskip}
\hypertarget{i.-main-improvements}{%
\subsection{I. Main Improvements}\label{i.-main-improvements}}

\Needspace{4\baselineskip}
\begin{enumerate}
\def\labelenumi{\arabic{enumi}.}
\tightlist
\item
  Value pretraining
\end{enumerate}

Before formal reinforcement learning (RL), train the critic (value model) specifically on a high-quality dataset. Early in RLHF, a randomly initialized critic makes extremely poor value estimates for long sequences. An incorrect critic misguides the actor and wastes computation. Pretraining lets the critic learn to recognize promising reasoning paths before ``starting work,'' so it can provide accurate long-term return estimates from the first RL step and avoid wasting substantial actor computation.

\Needspace{4\baselineskip}
\begin{enumerate}
\def\labelenumi{\arabic{enumi}.}
\setcounter{enumi}{1}
\tightlist
\item
  Decoupled GAE
\end{enumerate}

The bias--variance tradeoff is an important topic in reinforcement learning. Monte Carlo methods (discarding the critic network) are unbiased but have high variance and unstable training. In temporal-difference methods (using the critic network), the bias of each critic estimate continually accumulates, making it difficult for the crucial final-answer reward to propagate back to the first step in reasoning that spans hundreds of steps.

VAPO proposes decoupled GAE to combine the two.

\Needspace{5\baselineskip}
Its core advantage estimate can be written as:

\[
\hat{A}_t^{\mathrm{VAPO}}
=
\beta\cdot(R_{\mathrm{final}}-V(s_t))
+
(1-\beta)\cdot A_t^{\mathrm{GAE}}(\lambda).
\]

The first part is the GRPO/Monte Carlo component, using the actual final return to correct the direction; the second is the PPO/critic component, relying on temporal differences to provide smoother estimates.

\Needspace{5\baselineskip}
The actual return can also be forcibly injected into the GAE recursion:

\[
A_t
=
\delta_t
+
(\gamma\lambda)A_{t+1}
+
\eta\cdot(R-V(s_t)).
\]

The final term is the decoupled correction term, directly injecting the actual final return into the advantage estimate at each step.

\Needspace{5\baselineskip}
Here, \(R\) is \(G_t\) in the Monte Carlo estimate:

\[
G_t
=
r_t+\gamma r_{t+1}+\gamma^2 r_{t+2}
+\cdots+
\gamma^{T-t}R_{\mathrm{final}}.
\]

\begin{itemize}
\tightlist
\item
  In mathematical reasoning, rewards are usually provided only at the end, with \(\gamma=1\).
\item
  Therefore, \(G_t=R_{\mathrm{final}}\): it equals the final reward.
\end{itemize}

The most important role of decoupled GAE is to introduce accurate Monte Carlo estimates when the critic is biased early on, and use the critic to promote convergence when Monte Carlo becomes unstable later. Therefore, \(\beta\) generally decays over time, for example, from \(1.0\) to \(0.1\). (In fact, even without adjusting \(\beta\), the absolute value of \(R_{\mathrm{final}}-V(s_t)\) generally decreases as \(V(s_t)\) stabilizes, but decay is more effective.)

Early on, when the critic is weak, a stronger MC signal (the GRPO component) is needed to tell the model ``what is correct.'' The decoupled term ensures that even if the critic predicts \(V(s_{t+1})\) poorly, the final return \(R\) can bypass intermediate steps and correct the direction of \(A_t\) directly.

Later, when the critic is strong, it has learned fine-grained values at every step. At this point, \(A_t^{\mathrm{GAE}}\) (the PPO component) becomes more accurate and can distinguish subtleties such as ``the final answer is correct, but step 5 in the middle has a slight problem.''

\Needspace{5\baselineskip}
VAPO's advantage can be understood as:

\[
A_t^{\mathrm{VAPO}}
\approx
w_1\cdot A_t^{\mathrm{MC}}
+
w_2\cdot A_t^{\mathrm{Critic}}.
\]

\begin{itemize}
\tightlist
\item
  It uses GRPO's ``truth'' (actual returns) to correct critic bias.
\item
  It uses PPO's ``stability'' (the critic's dense signals) to reduce the variance of Monte Carlo sampling for long sequences.
\end{itemize}

This design enables VAPO to find the correct direction quickly early in large-model reasoning, like GRPO, and refine the logical quality of each step later, like PPO.

\Needspace{4\baselineskip}
\begin{enumerate}
\def\labelenumi{\arabic{enumi}.}
\setcounter{enumi}{2}
\tightlist
\item
  Length-adaptive GAE
\end{enumerate}

Like decoupled GAE, length-adaptive GAE is an important way to balance bias and variance. The difference is that decoupled GAE addresses ``whether the signal source is accurate'' (counteracting an initially incapable critic), while length-adaptive GAE addresses ``whether the signal propagates far enough'' (counteracting mathematical exponential decay), adapting to reasoning problems of different lengths.

\Needspace{5\baselineskip}
Define the effective horizon:

\[
\mathrm{Horizon}
\approx
\sum_{t=0}^{\infty}\lambda^t
=
\frac{1}{1-\lambda}.
\]

This is essentially defined using the expectation of a geometric distribution. It can be understood as continually receiving returns that decay by \(\lambda\) each time. In expectation, this is equivalent to a person starting at \(t\) whose sole job is to collect later returns to update \(A_t\). At each step, this person has probability \(\lambda\) of continuing forward and collecting the next return, and probability \(1-\lambda\) of being stopped (all subsequent returns are then ``intercepted'' and unavailable to \(A_t\)). In expectation, the person can reach round \(1/(1-\lambda)\), meaning that \(A_t\) can receive the first \(1/(1-\lambda)\) rounds of returns. This is the ``effective horizon.''

A fixed \(\lambda\) (such as \(0.95\), giving horizon \(=20\)) cannot adapt to different length requirements. For a 10-step answer, looking 20 steps ahead is too far (distant samples are weighted too heavily, and because they are obtained through Monte Carlo sampling, they are random and have high variance). For a 1000-step reasoning process, looking 20 steps ahead is extremely ``short-sighted'' (high bias, unable to see the terminal reward). We therefore want the effective horizon \(H\) to be proportional to sequence length \(L\) (\(H=\alpha L\)), obtaining \(\lambda=1-\frac{1}{\alpha L}\). The following uses \(\alpha=1\):

\begin{longtable}[]{@{}
  >{\raggedright\arraybackslash}p{(\columnwidth - 8\tabcolsep) * \real{0.1853}}
  >{\raggedleft\arraybackslash}p{(\columnwidth - 8\tabcolsep) * \real{0.24381}}
  >{\raggedright\arraybackslash}p{(\columnwidth - 8\tabcolsep) * \real{0.1853}}
  >{\raggedright\arraybackslash}p{(\columnwidth - 8\tabcolsep) * \real{0.1853}}
  >{\raggedright\arraybackslash}p{(\columnwidth - 8\tabcolsep) * \real{0.1853}}@{}}
\toprule
Sample type & Length \(L\) & Standard GAE (\(\lambda=0.95\)) & Adaptive GAE (\(\lambda=1-1/L\)) & Effect \\
\midrule
\endhead
Short dialogue & 20 & \(\lambda=0.95\) & \(\lambda=0.95\) & Both are identical, with a suitable horizon. \\
Medium-length reasoning & 100 & \(\lambda=0.95\) & \(\lambda=0.99\) & Standard GAE begins to lose sight of the ending; adaptive GAE raises \(\lambda\) to reduce decay. \\
Extremely long CoT & 1000 & \(\lambda=0.95\) & \(\lambda=0.999\) & Standard GAE's signal almost completely disappears; adaptive GAE pushes \(\lambda\) toward its limit, ensuring the reward at step 1000 reaches step 1. \\
\bottomrule
\end{longtable}

For long samples (long CoT), increasing \(\lambda\) toward 1 pushes the algorithm toward Monte Carlo behavior, tolerating high variance while ensuring the signal is not lost. For short samples (short QA), it retains a moderate \(\lambda\), using the low-variance properties of temporal differences for more stable training.

\Needspace{4\baselineskip}
\begin{enumerate}
\def\labelenumi{\arabic{enumi}.}
\setcounter{enumi}{3}
\tightlist
\item
  Token-level loss
\end{enumerate}

As in DAPO, average losses over all tokens rather than samples, distributing contributions to the final loss uniformly across tokens and avoiding dilution of crucial errors.

\Needspace{4\baselineskip}
\begin{enumerate}
\def\labelenumi{\arabic{enumi}.}
\setcounter{enumi}{4}
\tightlist
\item
  Group sampling
\end{enumerate}

Standard PPO assumes a stationary environment. In mathematical reasoning, however, if the model is weak, one sampled answer is likely to be incorrect (reward=0). If sampling continually produces incorrect answers, the critic learns nothing. VAPO's group sampling draws 4 or 8 answers per prompt at once. As long as one is correct, the critic can learn what is correct through comparison (using actual returns in decoupled GAE).

\Needspace{5\baselineskip}
\hypertarget{ii.-workflow-comparison-with-ppo}{%
\subsection{II. Workflow Comparison with PPO}\label{ii.-workflow-comparison-with-ppo}}

\Needspace{4\baselineskip}
\begin{enumerate}
\def\labelenumi{\arabic{enumi}.}
\tightlist
\item
  Initialization
\end{enumerate}

\Needspace{5\baselineskip}
Differences between PPO and VAPO initialization:

\begin{itemize}
\tightlist
\item
\Needspace{5\baselineskip}
  \textbf{PPO}:

  \begin{itemize}
  \tightlist
  \item
    Usually directly initializes the actor (from an SFT model) and critic (usually from a reward model).
  \item
    The distinction is that training starts directly, without a dedicated critic warm-up.
  \end{itemize}
\item
\Needspace{5\baselineskip}
  \textbf{VAPO}:

  \begin{itemize}
  \tightlist
  \item
    Additional step: value pretraining.
  \item
    Before entering the RL loop, the value model \(V_\psi\) is trained with supervision using SFT or reasoning data.
  \item
    The objective is to eliminate bias from the reward model so that the critic can predict long-sequence values relatively accurately on its first ``day at work,'' preventing early PPO collapse.
  \end{itemize}
\end{itemize}

\Needspace{4\baselineskip}
\begin{enumerate}
\def\labelenumi{\arabic{enumi}.}
\setcounter{enumi}{1}
\tightlist
\item
  Data sampling
\end{enumerate}

\begin{longtable}[]{@{}
  >{\raggedright\arraybackslash}p{(\columnwidth - 4\tabcolsep) * \real{0.33}}
  >{\raggedright\arraybackslash}p{(\columnwidth - 4\tabcolsep) * \real{0.33}}
  >{\raggedright\arraybackslash}p{(\columnwidth - 4\tabcolsep) * \real{0.33}}@{}}
\toprule
Dimension & Standard PPO & VAPO (augmented PPO) \\
\midrule
\endhead
Sampling strategy & Single/random sampling & Group sampling \\
Specific operation & For a prompt \(q\), usually sample 1 answer \(a\), or one answer for each different prompt in the batch. & For a prompt \(q\), force policy \(\pi_\theta\) to generate a group of \(G\) different answers \(\{o_1,\ldots,o_G\}\). \\
Data structure & Trajectories \(\{s_t,a_t,r_t\}\) are independent. & Data are organized in groups containing successful and unsuccessful samples. \\
Purpose & Obtain environmental feedback and calculate gradients. & Strengthen contrastive signals: multiple answers increase the probability of sampling a correct answer, ensuring surviving samples contain positive signals to learn from. \\
Reward calculation & Usually uses only outcome rewards or sparse rewards. & Dense rewards: usually combines outcome and process rewards. \\
Value calculation & Calculates \(V(s_t)\) during sampling. & Also calculates \(V(s_t)\), but for subsequent, more complex decoupled computation. \\
\bottomrule
\end{longtable}

Note: several prompts are processed simultaneously and stored in an experience replay buffer. For example, with 16 prompts, standard PPO generates one 64-token answer per prompt, storing 1024 data items in the buffer. VAPO generates 8 groups of 64-token answers per prompt, storing 8192 data items.

\Needspace{4\baselineskip}
\begin{enumerate}
\def\labelenumi{\arabic{enumi}.}
\setcounter{enumi}{2}
\tightlist
\item
  Advantage estimation
\end{enumerate}

\Needspace{5\baselineskip}
Differences between PPO and VAPO advantage estimation:

\Needspace{5\baselineskip}
\textbf{PPO}:

\begin{itemize}
\tightlist
\item
  The algorithm is standard generalized advantage estimation (GAE).
\item
\Needspace{5\baselineskip}
  Its formula is:
\end{itemize}

\[
A_t
=
\sum_{\ell\ge 0}(\gamma\lambda)^\ell\delta_{t+\ell}.
\]

\begin{itemize}
\tightlist
\item
  \(\lambda\) is a fixed constant, such as 0.95.
\item
  The limitation is that sequences of different lengths use the same horizon, easily losing signals in long sequences.
\end{itemize}

\Needspace{5\baselineskip}
\textbf{VAPO}:

\begin{itemize}
\tightlist
\item
  The algorithm combines decoupled GAE and length-adaptive GAE.
\item
  Step 1: calculate each answer \(o_i\)'s length \(l_i\).
\item
\Needspace{5\baselineskip}
  Step 2: calculate a separate \(\lambda_i\) for each sample, satisfying:
\end{itemize}

\[
\frac{1}{1-\lambda_i}
=
\alpha l_i.
\]

\begin{itemize}
\tightlist
\item
  Step 3: combine the critic's predictions \(V(s)\) and actual returns \(R\), using the decoupled formula to calculate \(\hat{A}_t\) and ensure signals penetrate the entire long sequence.
\end{itemize}

\Needspace{4\baselineskip}
\begin{enumerate}
\def\labelenumi{\arabic{enumi}.}
\setcounter{enumi}{3}
\tightlist
\item
  Optimization updates
\end{enumerate}

\begin{itemize}
\tightlist
\item
\Needspace{5\baselineskip}
  \textbf{PPO}:

  \begin{itemize}
  \tightlist
  \item
    Loss aggregation usually uses a sample-level mean, averaging the losses of all samples in a batch.
  \item
    The problem is that long texts (1000 tokens) and short texts (10 tokens) receive the same weight, diluting each long-text token's gradient by a factor of 100.
  \item
    Both \(\theta\) and \(\phi\) are updated.
  \end{itemize}
\item
\Needspace{5\baselineskip}
  \textbf{VAPO}:

  \begin{itemize}
  \tightlist
  \item
    Loss aggregation uses a token-level loss.
  \item
    Operationally, sum all token losses across all samples, \(\sum_i\sum_t \mathcal{L}_{i,t}\), then divide by the total token count.
  \item
    The purpose is to give tokens equal weight, ensuring every step in long-chain reasoning contributes sufficiently to the gradient and avoiding dilution of crucial errors.
  \end{itemize}
\end{itemize}

\FloatBarrier
\Needspace{5\baselineskip}
\hypertarget{references-110}{%
\subsection{References}\label{references-110}}

\vspace{0.25em}
\begin{itemize}
\tightlist
\item
  Yue, Y., Yuan, Y., Yu, Q., et al.~(2025). \href{https://arxiv.org/abs/2504.05118}{VAPO: Efficient and Reliable Reinforcement Learning for Advanced Reasoning Tasks}. arXiv:2504.05118.
\end{itemize}

\clearpage
